%% Title and Preprocessing
\documentclass[titlepage]{article}\pagestyle{empty}
\usepackage{amsmath}

\author{Kevin Fang}
\title{Homework \#7}
\date{\today}
\usepackage[left=1cm, right=2cm, vmargin=1cm]{geometry}
\usepackage{amsfonts}

\def\therefore{\boldsymbol{\text{ }
\leavevmode
\lower0.4ex\hbox{$\cdot$}
\kern-.5em\raise0.7ex\hbox{$\cdot$}
\kern-0.55em\lower0.4ex\hbox{$\cdot$}
\thinspace\text{ }}}

\usepackage{mathtools}
\DeclarePairedDelimiter\ceil{\lceil}{\rceil}
\DeclarePairedDelimiter\floor{\lfloor}{\rfloor}

% Start Document Here
\begin{document}
\maketitle

\pagebreak
\section*{Question 3.}
\subsection*{(a): 8.2.2}
\subsubsection*{B.} 

For $T(n) = n^3 + 3n^2 + 4$ and $h(n) = n^3$, $n^3 + 3n^2 + 4 = \Theta(n^3)$ if and only if $T(n) = O(h(n))$ and $T(n) = \Theta(h(n))$.\\
\begin{center}
For $n \geq 1$, $n \leq n^2 \leq n^3$.\\
\end{center}
\begin{center}
$n^3 + 3n^2 + 4 \leq n^3 + 3n^3 + 4n^3$\\
$n^3 + 3n^2 + 4 \leq 8n^3$\\
$n^3 + 3n^2 + 4 \leq 8h(n)$\\
\end{center}
$T(n) = O(h(n))$ since there are positive constants $c = 8$ and $n_0 = 1$ such that $T(n) \leq ch(n)$ when $n \geq n_0$.\\
\begin{center}
For $n \geq 1$, $3n^2 + 4 \geq 0$\\
\end{center}
\begin{center}
$3n^2 + 4 \geq 0$\\
$n^3 + 3n^2 + 4 \geq n^3$\\
$n^3 + 3n^2 + 4 \geq h(n)$\\
\end{center}
$T(n) = \Omega(h(n))$ since there are positive constants $c = 1$ and $n_0 = 1$ such that $T(n) \geq ch(n)$ when $n \geq n_0$.\\
$n^3 + 3n^2 + 4 = \Theta(n^3)$ since $T(n) = O(h(n))$ and $T(n) = \Theta(h(n))$.\\

\subsection*{(b): 8.3.5}
\subsubsection*{A.} 
This is a quick sort algorithm. The algorithm partitions the array around the pivot $p$. Specifically, it rearranges the array such that:\\
\begin{itemize}
\item All elements less than $p$ are moved to the left side of the array.
\item All elements greater than or equal to $p$ are moved to the right side of the array.
\end{itemize}
For example, if the input array is [3, -1, 2, 0, -2, 5] and p = 0, the output will be [-2, -1, 2, 0, 3, 5] or any other permutation where all elements less than 0 are on the left and all elements greater than or equal to 0 are on the right.\\
\subsubsection*{B.} 
The total number of times the lines $i:=i+1$ or $j:=j-1$ are executed depends on the distribution of numbers relative to $p$. Specifically:\\

\begin{itemize}
\item Maximizing Executions: If all elements are less than $p$, $i$ will be incremented $n-1$ times, and $j$ will not be decremented at all.
\item Minimizing Executions: If all elements are greater than or equal to $p$, $j$ will be decremented $n-1$ times, and $i$ will not be incremented at all.
\end{itemize}

The answer of the number of times the lines $i:=i+1$ or $j:=j-1$ are executed depends on the actual values of the numbers in the sequence, not just the length of the sequence. However, the total number of times the lines $i:=i+1$ or $j:=j-1$ are executed depends on the length of the sequence.\\
\subsubsection*{C.} 
The total number of times the swap operation is executed depends on how many elements are on the wrong side of $p$:

\begin{itemize}
\item Maximizing Swaps: If half of the elements are less than $p$ and half are greater than or equal to $p$, swaps will be maximized to $[\frac{n}{2}]$.
\item Minimizing Swaps: If all elements are either less than $p$ or greater than or equal to $p$, no swaps will occur.
\end{itemize}

The number of swaps depends on the actual values of the numbers in the sequence, not just the length of the sequence.
\subsubsection*{D.} The time complexity will be $\Omega(n)$ despite whether it is the worst-case or not.
\subsubsection*{E.} The time complexity will be $O(n)$.

\pagebreak
\section*{Question 4.}
\subsection*{(a): 5.1.2}
\subsubsection*{B.} $40^9 + 40^8 + 40^7$
\subsubsection*{C.} $14 * (40^8 + 40^7 + 40^6)$
\subsection*{(b): 5.3.2}
\subsubsection*{A.} $2^9 * 3$
\subsection*{(c): 5.3.3}
\subsubsection*{B.} $10 * 9 * 8 * 26^4$
\subsubsection*{C.} $10 * 9 * 8 * 26 * 25 * 24 * 23$
\subsection*{(d): 5.2.3}
\subsubsection*{A.}
We define a function $f:B^9 \rightarrow E_{10}$ by adding a 0 to the end of the string if it has even number of 1's and add a 1 to the end of the string if it has odd number of 1's.
The function is a bijection if since it has a well defined inverse $f^{-1}:E_{10}\rightarrow B^9$ that can be retrieved by removing the last bit.
\subsubsection*{B.} $\vert E_{10}\vert = \vert B_9 \vert=2^9=512$

\pagebreak
\section*{Question 5.}
\subsection*{(a): 5.4.2}
\subsubsection*{a.} $10^4 * 2=20000$
\subsubsection*{b.} $10 * 9 * 8 * 7 * 2=10080$
\subsection*{(b): 5.5.3}
\subsubsection*{a.} $2^{10}=1024$
\subsubsection*{b.} $2^7=128$
\subsubsection*{c.} $2^7 + 2^8=384$
\subsubsection*{d.} $2^8=256$
\subsubsection*{e.} $\binom{10}{6}=210$
\subsubsection*{f.} $\binom{9}{6}=84$
\subsubsection*{g.} $\binom{5}{1}\binom{5}{3}=50$
\subsection*{(c): 5.5.5}
\subsubsection*{a.} $\binom{30}{10}\binom{35}{10}=5515645706510940$
\subsection*{(d): 5.5.8}
\subsubsection*{c.} $\binom{26}{5}=65780$ 
\subsubsection*{d.} $\binom{13}{1}*\binom{48}{1}=624$
\subsubsection*{e.} $13 * 12 * \binom{4}{2} * \binom{4}{3} = 3744$
\subsubsection*{f.} $\binom{13}{5}*4^5=1317888$
\subsection*{(e): 5.6.6}
\subsubsection*{a.} $\binom{44}{5}\binom{56}{5}=4148350734528$
\subsubsection*{b.} $P(44, 2) * P(56, 2)=5827360$

\pagebreak
\section*{Question 6.}
\subsection*{(a): 5.7.2}
\subsubsection*{A.} $\binom{52}{5}-\binom{39}{5}=2023203$
\subsubsection*{B.} $\binom{52}{5}-(\binom{13}{5}*4^5)=1281072$
\subsection*{(b): 5.8.4}
\subsubsection*{A.} $5^{20}=95367431640625$
\subsubsection*{B.} $\binom{20}{4}\binom{16}{4}\binom{12}{4}\binom{8}{4}\binom{4}{4}=305540235000$

\pagebreak
\section*{Question 7.}
\subsection*{(a):} $0$
\subsection*{(b):} $P(5, 5)=120$
\subsection*{(c):} $P(6, 5)=720$
\subsection*{(d):} $P(7, 5)=2520$

\end{document}